\documentclass[10pt]{article}

\usepackage{graphicx}
\usepackage{tabularx}
\usepackage{fullpage}
\usepackage{caption}
\usepackage{cite}
\usepackage{url}
\usepackage{hyperref}
\usepackage{courier}

\usepackage{array}

\begin{document}
\title{Virtual slide viewer: feature development plan \vspace{-2em}}
\date{\today}
\maketitle

\section*{Subfeatures, functionality, or requirements}
\begin{center}
{
\renewcommand{\arraystretch}{1.3}
\begin{tabular}{| p{0.2\textwidth} | p{0.79\textwidth} |}
\hline
Basic layout & Separate tab with primary display pane for depiction of cells. \\
Cell shape & Cells rendered as circles, or ovals if bounding box is available. \\
Sample selector & One sample/slide viewed at a time, selectable e.g. with a dropdown menu. \\
Pan and zoom & Mouse pan and zoom. Continuous behavior is good but not necessary. \\
Downsampling & Reasonable downsampling when view is zoomed far out. \\
Phenotype selection & Legend with checkboxes for each phenotype to control colorization. \\
Phenotype representation & Cells colored according to selected phenotype (``composite phenotype''). \\
User-defined phenotypes & User-defined composite phenotypes (defined in existing other tab) also listed in phenotype checkboxes list. \\
Intensity imaging channel & Ability to overlay intensity (or 1/0 value) for one of the single imaged channels, e.g. as alpha channel or as color channel instead of (or in addition to?) selected phenotype colorization. \\
Lasso selection & Lasso selection of focus cell set, up to some reasonable maximum size (e.g. 5000 cells, since this set will be passed back and forth for stats summary purposes). \\
Context limitation & Lasso selection causes context limitation in feature matrix tab. \\ \hline
\end{tabular}
}
\end{center}

%  Mini-map... show selected/zoomed area in context?

\begin{itemize}
\itemsep0em
\item{Open SPT issues every time new data is needed from the API server, or the query formats need updating to accomodate this work. I (Jimmy) will work on these issues ASAP.}
\item{Real data not needed to complete prototype; can use random placement of e.g. 100k cells, and random assignment of e.g. 10 phenotypes, random assignment of a few channels of intensity.}
\end{itemize}

\section*{Schedule}

\begin{center}
{
\renewcommand{\arraystretch}{1.5}
\begin{tabular}{ | l | l | }
\hline
Date & Items completed \\ \hline
\begin{tabular}{ l r} January 26 (Friday) \\ (\emph{meeting 2:30pm EST}) \end{tabular} & Basic layout, Cell shape \\ \hline
\begin{tabular}{ l r} February 2 (Friday) \\ (\emph{meeting 2:30pm EST}) \end{tabular} & Sample selector, Pan and zoom, Downsampling \\ \hline
\begin{tabular}{ l r} February 9 (Friday)\end{tabular} & Phenotype selection, Phenotype representation, User-defined phenotypes \\ \hline
\begin{tabular}{ l r} February 16 (Friday) \\ (\emph{meeting 2:30pm EST}) \end{tabular} & Intensity imaging channel \\ \hline
\begin{tabular}{ l r} February 23 (Friday) \end{tabular} & Lasso selection, Context limitation \\ \hline
\begin{tabular}{ l r} February 26 (Monday) \\ (\emph{meeting 2:30pm EST}) \end{tabular} & \\ \hline
\begin{tabular}{ l r} March 1 (Friday) \end{tabular} & \emph{loose ends} \\ \hline
\end{tabular}
}
\end{center}

\end{document}
